\chapter{\texorpdfstring{$L^p$}{Lp} Spaces}

\begin{center}\itshape
A function is a point in a space, and two functions have a distance. Which distance? The largest gap, the average gap, the
energy of the gap? Each choice gives a space with its own geometry, and the right one depends on the question.
\end{center}

In 1888 an English mathematician, Leonard James Rogers, published an inequality about sums of products. A year later Otto
H\"older, in G\"ottingen, found it again, with a simpler proof, and history gave it his name. In 1906 Johan Jensen, who never
held a university post and worked all his life as an engineer for the Copenhagen Telephone Company, published the inequality
about convex functions that now carries his name. And in 1907, in the same volume of the \emph{Comptes rendus}, Frigyes Riesz
in Hungary and Ernst Fischer in Austria proved, independently, that the square-integrable functions form a complete space. Good
ideas, it seems, are in the air; but they still need someone to breathe them in.

What these people found is that functions can be treated as points, with a distance between them. The integral of $|f-g|$
measures how far apart two functions are on average; the integral of $|f-g|^2$ is an energy, and in probability a variance;
the essential supremum of $|f-g|$ is the worst case. The $L^p$ spaces put these distances in one family, indexed by
$1\le p\le\infty$, and three inequalities, Young's, H\"older's and Minkowski's, hold them together. Volume~1, Chapter~9, proved
H\"older's inequality for finite sums and announced that it was the key to $L^p$ spaces; here is the lock it opens.

The main theorem is that $L^p$ is complete (Riesz--Fischer): Cauchy sequences of functions have limits, which Riemann's
integral could never offer (Exercise~\ref{ex:L1incomplete}). Then come approximation, by step functions, by continuous
functions and by smooth ones, and the operation that does the smoothing: convolution, which is also what your phone does
when it blurs the background of a portrait.

\section{Young, H\"older, Minkowski}

\begin{lemma}[Young's inequality]\label{lem:young}
Let $1<p<\infty$ and $q=\frac p{p-1}$, so that $\frac1p+\frac1q=1$. For $a,b\ge0$,
\[ab\le\frac{a^p}p+\frac{b^q}q,\]
with equality if and only if $a^p=b^q$.
\end{lemma}
\begin{proof}
If $a=0$ or $b=0$ this is clear. Otherwise, since $\ln$ is strictly concave, $\ln(\frac1pa^p+\frac1qb^q)\ge\frac1p\ln a^p+\frac1q\ln b^q=\ln(ab)$, with equality iff $a^p=b^q$.
\end{proof}

The numbers $p$ and $q$ with $\frac1p+\frac1q=1$ are \emph{conjugate exponents}; $1$ and $\infty$ are conjugate too. Throughout, $(E,\mathcal A,\mu)$ is a measure space and, for a measurable $f$ and
$1\le p<\infty$,
\[\|f\|_p=\Bigl(\int|f|^p\,d\mu\Bigr)^{1/p}\in[0,+\infty].\]

\begin{theorem}[H\"older's inequality]\label{thm:holder}
Let $1<p<\infty$ and $q$ be conjugate. For measurable $f,g$, $\|fg\|_1\le\|f\|_p\,\|g\|_q$.
\end{theorem}
\begin{proof}
If $\|f\|_p=0$, then $f=0$ a.e. and both sides vanish; if $\|f\|_p=\infty$ and $\|g\|_q>0$ there is nothing to prove. So assume $0<\|f\|_p,\|g\|_q<\infty$ and, dividing $f$ and $g$ by their norms,
$\|f\|_p=\|g\|_q=1$. Integrating Young's inequality $|f||g|\le\frac{|f|^p}p+\frac{|g|^q}q$ gives $\|fg\|_1\le\frac1p+\frac1q=1$.
\end{proof}

For $p=q=2$ this is the Cauchy--Schwarz inequality, and with counting measure on $\{1,\dots,n\}$ it is the inequality of Volume~1.

\begin{theorem}[Minkowski's inequality]\label{thm:minkowski}
For $1\le p<\infty$ and measurable $f,g$, $\|f+g\|_p\le\|f\|_p+\|g\|_p$.
\end{theorem}
\begin{proof}
For $p=1$ integrate $|f+g|\le|f|+|g|$. Let $p>1$ and assume $\|f\|_p,\|g\|_p<\infty$. Then $\|f+g\|_p<\infty$, by convexity of $t\mapsto t^p$: $|f+g|^p\le2^{p-1}(|f|^p+|g|^p)$. Write
$|f+g|^p\le|f|\,|f+g|^{p-1}+|g|\,|f+g|^{p-1}$ and apply H\"older to each term, noting that $(p-1)q=p$:
\[\|f+g\|_p^p\le(\|f\|_p+\|g\|_p)\,\bigl\||f+g|^{p-1}\bigr\|_q=(\|f\|_p+\|g\|_p)\,\|f+g\|_p^{p-1}.\]
If $\|f+g\|_p>0$, divide by $\|f+g\|_p^{p-1}$.
\end{proof}

\section{The spaces \texorpdfstring{$L^p$}{Lp}}

Let $\mathcal L^p$ be the set of measurable $f$ (real or complex) with $\|f\|_p<\infty$. By Minkowski it is a vector space and $\|\cdot\|_p$ is a seminorm on it; $\|f\|_p=0$ exactly when $f=0$
a.e. So $\|\cdot\|_p$ is a norm on the quotient
\[L^p(\mu)=\mathcal L^p/\{f:f=0\text{ a.e.}\},\]
whose elements are classes of functions that agree almost everywhere. As is customary, we speak of ``the function $f\in L^p$'', meaning its class; nothing in integration theory can tell
two members of a class apart. For $p=\infty$, the \emph{essential supremum} is
\[\|f\|_\infty=\inf\{M\ge0:|f|\le M\text{ a.e.}\},\]
and the infimum is attained, since $\{|f|>\|f\|_\infty\}=\bigcup_n\{|f|>\|f\|_\infty+\frac1n\}$ is a countable union of null sets. Then $L^\infty(\mu)$ is the space of classes with $\|f\|_\infty<\infty$, a normed space
since $|f+g|\le\|f\|_\infty+\|g\|_\infty$ a.e. H\"older's inequality $\|fg\|_1\le\|f\|_1\|g\|_\infty$ is immediate. With counting measure on $\N$ these are the sequence spaces $\ell^p$.

\begin{theorem}[Riesz--Fischer]\label{thm:rieszfischer}
For $1\le p\le\infty$, $L^p(\mu)$ is a Banach space. If $f_n\to f$ in $L^p$ with $p<\infty$, a subsequence converges to $f$ almost everywhere.
\end{theorem}
\begin{proof}
Let $p<\infty$ and $(f_n)$ be Cauchy in $L^p$. Choose $n_1<n_2<\cdots$ with $\|f_{n_{k+1}}-f_{n_k}\|_p\le2^{-k}$, and let
\[G=|f_{n_1}|+\sum_{k\ge1}|f_{n_{k+1}}-f_{n_k}|.\]
The partial sums $G_K$ satisfy $\|G_K\|_p\le\|f_{n_1}\|_p+1$ by Minkowski, and $G_K^p\uparrow G^p$, so $\|G\|_p\le\|f_{n_1}\|_p+1$ by monotone convergence. Hence $G<\infty$ a.e. Where $G<\infty$, the series
$f_{n_1}+\sum_k(f_{n_{k+1}}-f_{n_k})$ converges absolutely; its partial sums are the $f_{n_K}$, so $f_{n_K}$ converges a.e. to a measurable $f$ (set $f=0$ elsewhere), and $|f|\le G$, $|f-f_{n_K}|\le2G$. By
dominated convergence with the integrable bound $(2G)^p$, $\|f-f_{n_K}\|_p\to0$. A Cauchy sequence with a convergent subsequence converges, so $f_n\to f$ in $L^p$. If $f_n\to f$ in $L^p$, then $(f_n)$ is Cauchy and
the argument provides a subsequence converging a.e., to a limit that must be $f$ a.e.

For $p=\infty$: the set $N=\bigcup_{m,n}\{|f_m-f_n|>\|f_m-f_n\|_\infty\}$ is null, and off $N$ the sequence is uniformly Cauchy, so it converges uniformly to a bounded $f$ (set $f=0$ on $N$); then
$\|f_n-f\|_\infty\le\sup_{x\notin N}|f_n(x)-f(x)|\to0$.
\end{proof}

\begin{think}
Can a sequence converge in $L^2([0,1])$ but at no point? And can it converge everywhere without converging in $L^2$?
\end{think}

Both happen. The typewriter of Exercise~\ref{ex:typewriter} has $\|f_n\|_2^2=\|f_n\|_1\to0$ and converges nowhere, while $\sqrt n\,\mathbf 1_{]0,\frac1n[}$ tends to $0$ everywhere and has $L^2$ norm $1$. So
the subsequence in Riesz--Fischer cannot be avoided, and convergence in $L^p$ and almost everywhere are not comparable in general; Vitali's theorem (Theorem~\ref{thm:vitali7}) describes the link for $p=1$.

\emph{Inclusions.} If $\mu(E)<\infty$ and $1\le p<q\le\infty$, H\"older with exponent $\frac qp$ applied to $|f|^p\cdot1$ gives
\[\|f\|_p\le\mu(E)^{\frac1p-\frac1q}\|f\|_q,\]
so $L^q\subseteq L^p$: on a finite measure space, a larger exponent means a smaller space. For sequences it is the other way round: if $\|a\|_p\le1$, then $|a_n|\le1$, so $|a_n|^q\le|a_n|^p$ and
$\|a\|_q\le1$; by homogeneity $\|a\|_q\le\|a\|_p$ and $\ell^p\subseteq\ell^q$. On $\R$ neither inclusion holds: large exponents forbid singularities, small exponents forbid slowly decaying tails.

\begin{think}
For which $p$ is $x\mapsto\frac1{\sqrt x}$ in $L^p(]0,1[)$? In $L^p(]1,\infty[)$? In $L^p(]0,\infty[)$?
\end{think}

The integral $\int_0^1x^{-p/2}\,dx$ is finite iff $\frac p2<1$, and $\int_1^\infty x^{-p/2}\,dx<\infty$ iff $\frac p2>1$. So $\frac1{\sqrt x}$ is in $L^p(]0,1[)$ for $p<2$, in $L^p(]1,\infty[)$ for $p>2$ (including $p=\infty$), and in no $L^p(]0,\infty[)$:
the singularity at $0$ and the tail at $\infty$ ask for opposite exponents, and $p=2$ fails at both, logarithmically.

\begin{figure}[ht]
\centering
\fig{ch09-balls}
\caption{The unit balls $\{x\in\R^2:\|x\|_p\le1\}$ for $p=1$, $1.5$, $2$, $4$ and $\infty$. They increase with $p$, as $\ell^p\subseteq\ell^q$ for $p<q$ predicts. For $1<p<\infty$ they are strictly convex
and smooth; $p=1$ and $p=\infty$ have corners and flat sides, a difference that will matter in Volume~8.}
\label{fig:lpballs}
\end{figure}

\begin{theorem}[Jensen's inequality]\label{thm:jensen}
Let $\mu$ be a probability, $I$ an open interval, $\varphi:I\to\R$ convex and $f\in L^1(\mu)$ with values in $I$. Then $\int f\,d\mu\in I$, the negative part of $\varphi\circ f$ is integrable, and
\[\varphi\Bigl(\int f\,d\mu\Bigr)\le\int\varphi\circ f\,d\mu.\]
\end{theorem}
\begin{proof}
Let $m=\int f\,d\mu$. If $I$ has a finite left end $a$, then $f-a>0$, so $m-a=\int(f-a)>0$; similarly on the right: $m\in I$. For $s<m<t$ in $I$, convexity gives the three-slope inequality
$\frac{\varphi(m)-\varphi(s)}{m-s}\le\frac{\varphi(t)-\varphi(m)}{t-m}$ (write $m$ as a convex combination of $s$ and $t$). Let $c$ be the supremum of the left-hand sides over $s<m$; then
$\varphi(t)\ge\varphi(m)+c(t-m)$ for $t>m$, and also for $t<m$ (as $\frac{\varphi(m)-\varphi(t)}{m-t}\le c$). So $\varphi$ lies above the line $\ell(t)=\varphi(m)+c(t-m)$. In particular $\varphi$ is bounded below by
affine functions near each point and, being convex, continuous on $I$ (the slopes are bounded on compact subintervals), so $\varphi\circ f$ is measurable. Now $\varphi\circ f\ge\ell\circ f$, an integrable function, so
$(\varphi\circ f)^-$ is integrable, and $\int\varphi\circ f\ge\int\ell\circ f=\varphi(m)+c(m-m)=\varphi(m)$.
\end{proof}

With $\varphi(t)=t^{q/p}$ and $|f|^p$ in place of $f$, Jensen gives $\|f\|_p\le\|f\|_q$ for $p<q$ on a probability space, the inclusion above with $\mu(E)=1$. With $\varphi=\exp$ and counting
probabilities it gives the inequality between arithmetic and geometric means (Exercise~\ref{ex:jensen}). In probability it reads $\varphi(\mathbb E[X])\le\mathbb E[\varphi(X)]$, and Volume~11 extends it to conditional expectations.

\section{Approximation}

\begin{proposition}\label{prop:simpledense}
For $1\le p\le\infty$, the simple functions in $L^p$ are dense in $L^p$. For $p<\infty$ these are the simple functions $\sum_ia_i\mathbf 1_{A_i}$ with $\mu(A_i)<\infty$.
\end{proposition}
\begin{proof}
For $f\ge0$ in $L^p$, take simple $0\le s_n\uparrow f$ (Proposition~\ref{prop:measurablefns}). If $p<\infty$, $|f-s_n|^p\le f^p$, integrable, so $\|f-s_n\|_p\to0$ by dominated convergence; and $s_n^p\le f^p$ forces
$\mu(\{s_n\ge a\})<\infty$ for $a>0$. If $p=\infty$, changing $f$ on a null set makes it bounded, and the $s_n$ of Proposition~\ref{prop:measurablefns} converge uniformly on $\{f\le M\}$. A general $f$ is a combination
of four non-negative ones.
\end{proof}

On $\R^d$, a \emph{step function} is a finite linear combination of indicators of boxes, and $C_c(\R^d)$ is the space of continuous functions with compact support.

\begin{theorem}\label{thm:ccdense}
For $1\le p<\infty$, step functions and functions of $C_c(\R^d)$ are dense in $L^p(\R^d)$. Neither is dense in $L^\infty(\R^d)$.
\end{theorem}
\begin{proof}
By Proposition~\ref{prop:simpledense} it suffices to approximate $\mathbf 1_A$ with $\lambda_d(A)<\infty$. By regularity (Proposition~\ref{prop:regularity}) there are a compact $K$ and an open $U$ with $K\subseteq A\subseteq U$
and $\lambda_d(U\setminus K)<\varepsilon$; intersecting $U$ with a large open ball, we may take $U$ bounded.

\emph{Continuous.} The Urysohn function $g(x)=\frac{d(x,U^c)}{d(x,U^c)+d(x,K)}$ (Proposition~\ref{prop:urysohn}) is continuous, equal to $1$ on $K$ and $0$ off $U$, with values in $[0,1]$; its support lies in $\overline U$, which is compact.
Then $|\mathbf 1_A-g|\le\mathbf 1_{U\setminus K}$, so $\|\mathbf 1_A-g\|_p^p\le\varepsilon$.

\emph{Steps.} $U$ is a countable union of closed dyadic cubes $Q_k$ with disjoint interiors (as in the proof of Theorem~\ref{thm:cov}), and $\sum_k\lambda_d(Q_k)=\lambda_d(U)<\infty$. For $N$ large,
$\sum_{k>N}\lambda_d(Q_k)<\varepsilon$, and $h=\mathbf 1_{Q_1\cup\cdots\cup Q_N}$, a step function (up to the null boundaries), satisfies $|\mathbf 1_A-h|\le\mathbf 1_{U\setminus K}+\mathbf 1_{\bigcup_{k>N}Q_k}$, so
$\|\mathbf 1_A-h\|_p^p\le2\varepsilon$.

For $L^\infty(\R)$, let $f=\mathbf 1_{[0,\infty[}$ and $g$ be continuous with $\|f-g\|_\infty=c<\frac12$. Then $g\ge1-c$ a.e. on $]0,1[$, hence everywhere there (the set where it holds is dense and $g$ is
continuous), and $g(0)\ge1-c$; likewise $g\le c$ on $]-1,0[$, so $g(0)\le c<1-c$, a contradiction. So $C_c$ is at distance at least $\frac12$ from $f$. A step function $h$ vanishes outside a bounded set,
so $\|f-h\|_\infty\ge1$.
\end{proof}

\begin{proposition}[continuity of translation]\label{prop:translation}
For $f\in L^p(\R^d)$ with $1\le p<\infty$ and $\tau_hf(x)=f(x-h)$, $\|\tau_hf-f\|_p\to0$ as $h\to0$. This fails in $L^\infty$.
\end{proposition}
\begin{proof}
For $g\in C_c$, supported in a ball $B$, $g$ is uniformly continuous, so $\sup|\tau_hg-g|\to0$, and for $|h|\le1$ the support of $\tau_hg-g$ lies in the $1$-neighbourhood $B'$ of $B$: $\|\tau_hg-g\|_p\le\lambda_d(B')^{1/p}\sup|\tau_hg-g|\to0$.
For $f\in L^p$ and $\varepsilon>0$, pick $g\in C_c$ with $\|f-g\|_p<\varepsilon$ (Theorem~\ref{thm:ccdense}); translation invariance of $\lambda_d$ gives $\|\tau_hf-\tau_hg\|_p=\|f-g\|_p$, so
$\|\tau_hf-f\|_p\le2\varepsilon+\|\tau_hg-g\|_p<3\varepsilon$ for $h$ small. In $L^\infty$, $\|\tau_h\mathbf 1_{[0,\infty[}-\mathbf 1_{[0,\infty[}\|_\infty=1$ for every $h\ne0$.
\end{proof}

\section{Convolution}

For measurable $f,g$ on $\R^d$, the \emph{convolution} is
\[(f*g)(x)=\int_{\R^d}f(x-y)\,g(y)\,dy,\]
wherever the integral converges absolutely. Every class in $L^p$ contains a Borel function (Exercise~\ref{ex:completion} of Chapter~7 describes the Lebesgue measurable sets as Borel sets up to null sets), and we always use
Borel representatives, so that $(x,y)\mapsto f(x-y)g(y)$ is Borel on $\R^{2d}$ and Tonelli applies. Changing $f$ or $g$ on a null set does not change $f*g$, and the substitution $y\mapsto x-y$ gives $f*g=g*f$.
Chapter~9 showed that if $f$ and $g$ are probability densities, $f*g$ is the density of the sum of independent variables with these densities.

\begin{proposition}\label{prop:convolution}
\begin{enumerate}[label=(\alph*)]
\item If $f\in L^1$ and $g\in L^p$, $1\le p\le\infty$, then $f*g(x)$ is defined for almost every $x$, and $\|f*g\|_p\le\|f\|_1\|g\|_p$.
\item If $f\in L^p$ and $g\in L^q$ with $p,q$ conjugate, then $f*g$ is defined everywhere, bounded by $\|f\|_p\|g\|_q$, and uniformly continuous.
\end{enumerate}
\end{proposition}
\begin{proof}
(a) For $p=\infty$, $|f*g(x)|\le\|g\|_\infty\int|f(x-y)|\,dy=\|f\|_1\|g\|_\infty$. For $p<\infty$, H\"older's inequality for the measure $|f(x-y)|\,dy$, of total mass $\|f\|_1$, gives (Jensen's form for $p>1$, trivial
for $p=1$)
\[\Bigl(\int|f(x-y)|\,|g(y)|\,dy\Bigr)^p\le\|f\|_1^{p-1}\int|f(x-y)|\,|g(y)|^p\,dy.\]
Integrating in $x$ with Tonelli, $\int(\int|f(x-y)||g(y)|\,dy)^pdx\le\|f\|_1^{p-1}\cdot\|f\|_1\|g\|_p^p$. So the integral defining $f*g(x)$ converges absolutely for a.e. $x$, and $\|f*g\|_p\le\|f\|_1\|g\|_p$.
(b) By H\"older and translation invariance, $|f*g(x)|\le\|\tau_xf^\vee\|_p\|g\|_q=\|f\|_p\|g\|_q$, where $f^\vee(y)=f(-y)$. One of $p,q$ is finite, say $p$ (otherwise exchange $f$ and $g$); then
$|f*g(x+h)-f*g(x)|\le\|\tau_{-h}f-f\|_p\|g\|_q\to0$ uniformly in $x$, by Proposition~\ref{prop:translation}.
\end{proof}

Exercise~\ref{ex:youngconv} proves the general \emph{Young inequality} $\|f*g\|_r\le\|f\|_p\|g\|_q$ for $\frac1p+\frac1q=1+\frac1r$, which contains both (a) and (b).

\emph{Approximate identities.} Let $\varphi\ge0$ be integrable with $\int\varphi=1$, and $\varphi_t(x)=t^{-d}\varphi(\frac xt)$ for $t>0$. Then $\int\varphi_t=1$, and as $t\to0$ the mass of $\varphi_t$ concentrates at $0$:
$\int_{|x|>\delta}\varphi_t=\int_{|y|>\delta/t}\varphi\to0$. Convolving with $\varphi_t$ averages $f$ over a region of size about $t$.

\begin{theorem}[approximate identities]\label{thm:approxid}
Let $\varphi_t$ be as above.
\begin{enumerate}[label=(\alph*)]
\item If $f\in L^p(\R^d)$, $1\le p<\infty$, then $\|f*\varphi_t-f\|_p\to0$ as $t\to0$.
\item If $f$ is bounded and measurable, then $f*\varphi_t(x)\to f(x)$ at every point $x$ where $f$ is continuous, uniformly on any set where $f$ is uniformly continuous.
\end{enumerate}
\end{theorem}
\begin{proof}
Substituting $y=tz$, $f*\varphi_t(x)-f(x)=\int(f(x-tz)-f(x))\,\varphi(z)\,dz$.
(a) Since $\varphi(z)\,dz$ is a probability, Jensen's inequality with $t\mapsto|t|^p$ (or H\"older) gives $|f*\varphi_t(x)-f(x)|^p\le\int|f(x-tz)-f(x)|^p\varphi(z)\,dz$. Integrating in $x$ and using Tonelli,
\[\|f*\varphi_t-f\|_p^p\le\int\|\tau_{tz}f-f\|_p^p\,\varphi(z)\,dz.\]
For each $z$ the integrand tends to $0$ (Proposition~\ref{prop:translation}), and it is bounded by $2^p\|f\|_p^p\varphi(z)$, integrable: dominated convergence.
(b) Let $|f|\le M$, $\varepsilon>0$, and $\delta>0$ with $|f(x-y)-f(x)|<\varepsilon$ for $|y|<\delta$. Then
\[|f*\varphi_t(x)-f(x)|\le\varepsilon\int\varphi_t+2M\int_{|y|\ge\delta}\varphi_t(y)\,dy,\]
and the last integral tends to $0$. The same $\delta$ works for all $x$ in a set where $f$ is uniformly continuous.
\end{proof}

\begin{figure}[ht]
\centering
\fig{ch09-mollify}
\caption{The indicator of $[-1,1]$ convolved with the mollifiers $\varphi_t$ for $t=0.5$, $0.25$ and $0.1$. Each convolution is smooth; it equals $1$ and $0$ at distance more than $t$ from the
jumps, and converges to the step in $L^1$ and at every point except $\pm1$, where it tends to $\frac12$.}
\label{fig:mollify}
\end{figure}

\emph{Mollifiers.} The function $\theta(s)=e^{-1/s}$ for $s>0$, $\theta(s)=0$ for $s\le0$, is $C^\infty$ on $\R$: on $]0,\infty[$, induction shows $\theta^{(k)}(s)=P_k(\frac1s)e^{-1/s}$ for a polynomial $P_k$, which tends to $0$
as $s\to0^+$ since $e^{-u}$ beats every power of $u$ at infinity; so each derivative extends continuously by $0$, and the mean value theorem shows that the extension is differentiable at $0$
with derivative $0$. Hence $\varphi(x)=c\,\theta(1-|x|^2)$, with $c$ chosen so that $\int\varphi=1$, is $C^\infty$, non-negative, and supported in the closed unit ball. Such a $\varphi$ is a \emph{mollifier}.

\begin{theorem}\label{thm:smoothdense}
If $f\in L^p(\R^d)$, $1\le p\le\infty$, and $\varphi$ is a mollifier, then $f*\varphi_t$ is $C^\infty$, with $\partial^\alpha(f*\varphi_t)=f*\partial^\alpha\varphi_t$. For $p<\infty$, the space $C_c^\infty(\R^d)$ of smooth functions
with compact support is dense in $L^p(\R^d)$.
\end{theorem}
\begin{proof}
Fix $x_0$ and consider $x$ with $|x-x_0|<1$. The functions $y\mapsto f(y)\,\partial_i\varphi_t(x-y)$ are dominated by $\sup|\partial_i\varphi_t|\cdot|f(y)|\mathbf 1_{|y-x_0|\le1+t}$, which is integrable because $f\in L^p$ is
integrable on bounded sets (by the inclusion for finite measure). Differentiating under the integral sign (Theorem~\ref{thm:parameter}(b), one variable at a time) gives
$\partial_i(f*\varphi_t)=f*\partial_i\varphi_t$, which is continuous by the same domination and Theorem~\ref{thm:parameter}(a); induct on the order. For density, given $f\in L^p$ and $\varepsilon>0$, pick $g\in C_c$ with
$\|f-g\|_p<\varepsilon$ (Theorem~\ref{thm:ccdense}). Then $g*\varphi_t$ is smooth, has compact support (within $t$ of that of $g$), and $\|g*\varphi_t-g\|_p<\varepsilon$ for small $t$, by Theorem~\ref{thm:approxid}.
\end{proof}

Convolution thus turns any $L^p$ function into a smooth one, as close as we wish. This is how one proves identities for all of $L^p$: check them on smooth compactly supported functions, where
calculus is available, and pass to the limit by continuity. Volume~8 uses it constantly, for Fourier transforms, distributions and Sobolev spaces.

\begin{beauty}
Completeness is what turns analysis into geometry. Since $L^2$ is complete and its norm comes from the inner product $\int f\bar g$, it is a Hilbert space, where orthogonal projections exist and
Pythagoras holds for functions: that is the starting point of Volume~8. In probability, $L^2$ is the space of random variables of finite variance, and the variance is a squared distance. A
single inequality about two numbers, $ab\le\frac{a^p}p+\frac{b^q}q$, holds all of it together.
\end{beauty}

\section{Exercises}

\begin{exercise}[$\star\star\star$]\label{ex:youngconv}
\emph{Young's inequality for convolutions.} (a) Prove H\"older's inequality for three functions: $\|fgh\|_1\le\|f\|_a\|g\|_b\|h\|_c$ when $\frac1a+\frac1b+\frac1c=1$. (b) Let $1\le p,q,r\le\infty$ with
$\frac1p+\frac1q=1+\frac1r$. Show that $\|f*g\|_r\le\|f\|_p\|g\|_q$, treating first the case $1<p,q<r<\infty$ by writing $|f(x-y)g(y)|$ as a product of three factors.
\end{exercise}

\begin{exercise}[$\star\star$]\label{ex:approxid}
\emph{Heat and Poisson kernels.} (a) Show that $G_t(x)=\frac1{\sqrt{4\pi t}}e^{-x^2/(4t)}$ and $P_t(x)=\frac t{\pi(t^2+x^2)}$ are approximate identities on $\R$ as $t\to0$, of the form $\varphi_s$ for suitable $\varphi$ and
$s$. (b) For $f\in L^p(\R)$, show that $u(t,x)=(f*G_t)(x)$ solves the heat equation $\partial_tu=\partial_x^2u$ for $t>0$ and that $u(t,\cdot)\to f$ in $L^p$ as $t\to0$. (c) Show that $G_s*G_t=G_{s+t}$ and
$P_s*P_t=P_{s+t}$, using Chapter~9 and the laws of sums of independent Gaussian, respectively Cauchy, variables (for Cauchy, admit the characteristic function $e^{-t|\xi|}$, proved in Volume~8).
\end{exercise}

\begin{exercise}[$\star\star$]\label{ex:logconvex}
\emph{Interpolation in miniature.} Let $1\le p_0<p_1\le\infty$, $0<\theta<1$, and $\frac1{p_\theta}=\frac{1-\theta}{p_0}+\frac\theta{p_1}$. Show that $\|f\|_{p_\theta}\le\|f\|_{p_0}^{1-\theta}\|f\|_{p_1}^\theta$. Deduce that
$L^{p_0}\cap L^{p_1}\subseteq L^p$ for all $p$ between them, and that $\frac1p\mapsto\ln\|f\|_p$ is convex.
\end{exercise}

\begin{exercise}[$\star\star$]\label{ex:pinfty}
Let $f\in L^{p_0}\cap L^\infty$ for some $p_0<\infty$. Show that $\|f\|_p\to\|f\|_\infty$ as $p\to\infty$. Give an example where $f\in L^\infty$, $f\notin L^p$ for every $p<\infty$.
\end{exercise}

\begin{exercise}[$\star\star\star$]\label{ex:hardy}
\emph{Hardy's inequality.} Let $p>1$, $f\ge0$ in $L^p(]0,\infty[)$, and $F(x)=\frac1x\int_0^xf(s)\,ds$. Show that $\|F\|_p\le\frac p{p-1}\|f\|_p$: write $F(x)=\int_0^1f(xu)\,du$, insert $u^\alpha u^{-\alpha}$ and use
H\"older in $u$, then integrate in $x$ and choose $\alpha$ well. Show that the constant $\frac p{p-1}$ cannot be improved, with $f(x)=x^{-1/p}\mathbf 1_{[1,N]}(x)$.
\end{exercise}

\begin{exercise}[$\star$]\label{ex:jensen}
Use Jensen's inequality to prove: (a) the inequality of arithmetic and geometric means $(x_1\cdots x_n)^{1/n}\le\frac{x_1+\cdots+x_n}n$ for $x_i>0$; (b) $\mathbb E[X]^2\le\mathbb E[X^2]$; (c) $\int\ln f\,d\mu\le\ln\int f\,d\mu$ for
a probability $\mu$ and $f>0$ integrable.
\end{exercise}

\begin{exercise}[$\star\star$]\label{ex:linfinity}
(a) Show that $L^\infty([0,1])$ is not separable: find uncountably many functions at mutual distance $1$. (b) Show that $L^p([0,1])$ is separable for $p<\infty$. (c) Show that $\sin$ is at distance at
least $1$, in $L^\infty(\R)$, from every step function, and at distance at least $1$ from every function of $C_c(\R)$.
\end{exercise}

\section{Hints}

\begin{hint}{ex:youngconv}
(a) Apply H\"older twice, first with exponents $a$ and $\frac a{a-1}$. (b) With $\frac1{s_1}=\frac1p-\frac1r$ and $\frac1{s_2}=\frac1q-\frac1r$, write $|f(x-y)g(y)|=(|f(x-y)|^p|g(y)|^q)^{1/r}\cdot|f(x-y)|^{p/s_1}\cdot|g(y)|^{q/s_2}$.
\end{hint}
\begin{hint}{ex:approxid}
(a) $G_t=\varphi_{\sqrt t}$ with $\varphi=G_1$, and $P_t=\varphi_t$ with $\varphi=P_1$. (b) Check $\partial_tG_t=\partial_x^2G_t$ and differentiate under the integral sign on $t\ge t_0>0$. (c) $G_t$ is the density of
$\mathcal N(0,2t)$.
\end{hint}
\begin{hint}{ex:logconvex}
Write $|f|^{p_\theta}=|f|^{(1-\theta)p_\theta}|f|^{\theta p_\theta}$ and apply H\"older with exponents $\frac{p_0}{(1-\theta)p_\theta}$ and $\frac{p_1}{\theta p_\theta}$.
\end{hint}
\begin{hint}{ex:pinfty}
Upper bound: $|f|^p\le\|f\|_\infty^{p-p_0}|f|^{p_0}$. Lower bound: for $M<\|f\|_\infty$ the set $\{|f|>M\}$ has positive, finite measure.
\end{hint}
\begin{hint}{ex:hardy}
The choice $\alpha=\frac1{pq}$, with $q$ conjugate to $p$, makes the two integrals in $u$ balance. For the sharpness, compute $F$ for $x\ge N$ and for $1\le x\le N$.
\end{hint}
\begin{hint}{ex:jensen}
(a) $\exp$ is convex; use the uniform probability on $\{1,\dots,n\}$ and $f(i)=\ln x_i$.
\end{hint}
\begin{hint}{ex:linfinity}
(a) $\mathbf 1_{[0,t]}$ for $t\in[0,1]$. (b) Step functions with rational heights and rational endpoints. (c) Step functions and functions of $C_c$ vanish outside a bounded set.
\end{hint}

\section{Solutions}

\begin{solution}{ex:youngconv}
\idea Young's inequality for convolutions interpolates between three easy cases: $\|f*g\|_\infty\le\|f\|_p\|g\|_q$ (H\"older), $\|f*g\|_p\le\|f\|_1\|g\|_p$ and $\|f*g\|_1\le\|f\|_1\|g\|_1$. The trick is to split the integrand
into three factors, one of which carries the whole $L^r$ integrability, and the other two of which are controlled uniformly in $x$.
\writeup (a) H\"older with exponents $a$ and $a'=\frac a{a-1}$ gives $\|fgh\|_1\le\|f\|_a\|gh\|_{a'}$, and $\|gh\|_{a'}^{a'}=\int|g|^{a'}|h|^{a'}\le\||g|^{a'}\|_{b/a'}\||h|^{a'}\|_{c/a'}=\|g\|_b^{a'}\|h\|_c^{a'}$, since $\frac{a'}b+\frac{a'}c=a'(1-\frac1a)=1$.
(b) The case $r=\infty$ is part (b) of Proposition~\ref{prop:convolution}; $p=1$ forces $q=r$ and $q=1$ forces $p=r$, which are part (a); so assume $1<p,q<r<\infty$. Then $s_1,s_2$ defined in the hint are in
$]1,\infty[$, and $\frac1r+\frac1{s_1}+\frac1{s_2}=\frac1p+\frac1q-\frac1r=1$. Since $\frac pr+\frac p{s_1}=1$ and $\frac qr+\frac q{s_2}=1$, the factorisation of the hint holds, and (a) with exponents $r,s_1,s_2$ gives
\[\int|f(x-y)g(y)|\,dy\le\Bigl(\int|f(x-y)|^p|g(y)|^qdy\Bigr)^{1/r}\|f\|_p^{p/s_1}\|g\|_q^{q/s_2}.\]
Raise to the power $r$ and integrate in $x$ with Tonelli: $\int(\ldots)^rdx\le\|f\|_p^p\|g\|_q^q\cdot\|f\|_p^{rp/s_1}\|g\|_q^{rq/s_2}=\|f\|_p^r\|g\|_q^r$, because $p+\frac{rp}{s_1}=p+r-p=r$, and likewise for $q$. So $f*g(x)$ is defined
a.e. and $\|f*g\|_r\le\|f\|_p\|g\|_q$.
\end{solution}

\begin{solution}{ex:approxid}
\idea Both kernels are dilations of one fixed probability density, so Theorem~\ref{thm:approxid} applies directly. The heat equation holds for the kernel itself, and differentiation under the
integral sign carries it to $f*G_t$. The semigroup identities are statements about sums of independent variables.
\writeup (a) $G_1(x)=\frac{e^{-x^2/4}}{\sqrt{4\pi}}$ is the $\mathcal N(0,2)$ density, and $s^{-1}G_1(\frac xs)=\frac1{s\sqrt{4\pi}}e^{-x^2/(4s^2)}=G_{s^2}(x)$, so $G_t=(G_1)_{\sqrt t}$. $P_1(x)=\frac1{\pi(1+x^2)}$ has integral
$\frac1\pi[\arctan x]_{-\infty}^\infty=1$ and $t^{-1}P_1(\frac xt)=P_t(x)$. As $t\to0$, $\sqrt t\to0$, so both are approximate identities. (b) A direct computation gives
$\partial_tG_t(x)=G_t(x)(\frac{x^2}{4t^2}-\frac1{2t})=\partial_x^2G_t(x)$. For $t\in[t_0,t_1]$ with $t_0>0$, the derivatives $\partial_tG_t$, $\partial_xG_t$, $\partial^2_xG_t$ are bounded by $C(1+x^2)e^{-x^2/(4t_1)}$, a function in every $L^q$;
for $|x|\le R$ these bounds at $x-y$ are at most $k(y)=C(1+(|y|+R)^2)e^{-((|y|-R)^+)^2/(4t_1)}$, which is in $L^q$, so $|f|k$ is integrable by H\"older and Theorem~\ref{thm:parameter} allows differentiation under the integral: $\partial_tu=f*\partial_tG_t=f*\partial_x^2G_t=\partial_x^2u$.
By (a) and Theorem~\ref{thm:approxid}, $u(t,\cdot)\to f$ in $L^p$. (c) $G_s*G_t$ is the density of the sum of independent $\mathcal N(0,2s)$ and $\mathcal N(0,2t)$ variables, which is $\mathcal N(0,2(s+t))$ (Volume~4),
so $G_s*G_t=G_{s+t}$. $P_t$ is the density of $tC$ with $C$ standard Cauchy, whose characteristic function is $e^{-t|\xi|}$; the sum of independent $sC$ and $tC'$ has characteristic function
$e^{-(s+t)|\xi|}$, so it has the law of $(s+t)C$ (a law is determined by its characteristic function, Volume~4), and $P_s*P_t=P_{s+t}$. The heat kernel and the Poisson kernel thus form \emph{semigroups}:
diffusing for a time $s$, then $t$, is diffusing for $s+t$.
\end{solution}

\begin{solution}{ex:logconvex}
\idea The exponent $p_\theta$ is an average of $p_0$ and $p_1$ in the scale $\frac1p$. Splitting $|f|^{p_\theta}$ into two powers, one measured in $L^{p_0}$ and one in $L^{p_1}$, is exactly what H\"older can do with the
right exponents.
\writeup Let $p_1<\infty$ first. Put $a=\frac{p_0}{(1-\theta)p_\theta}$ and $b=\frac{p_1}{\theta p_\theta}$; then $\frac1a+\frac1b=p_\theta(\frac{1-\theta}{p_0}+\frac\theta{p_1})=1$. H\"older gives
\[\int|f|^{p_\theta}\le\Bigl(\int|f|^{p_0}\Bigr)^{1/a}\Bigl(\int|f|^{p_1}\Bigr)^{1/b}=\|f\|_{p_0}^{(1-\theta)p_\theta}\|f\|_{p_1}^{\theta p_\theta},\]
and taking the power $\frac1{p_\theta}$ gives the claim. If $p_1=\infty$, then $p_\theta=\frac{p_0}{1-\theta}$ and $\int|f|^{p_\theta}\le\|f\|_\infty^{p_\theta-p_0}\int|f|^{p_0}$, which is the same inequality. Every $p$ between $p_0$ and $p_1$ is a
$p_\theta$, so $L^{p_0}\cap L^{p_1}\subseteq L^p$. Finally, with $s=\frac1p$, the inequality says $\ln\|f\|_{1/s}\le(1-\theta)\ln\|f\|_{1/s_0}+\theta\ln\|f\|_{1/s_1}$ for $s=(1-\theta)s_0+\theta s_1$: convexity in $s$. This is the
shadow, for a single function, of the Riesz--Thorin theorem for operators (Volume~8).
\end{solution}

\begin{solution}{ex:pinfty}
\idea $\|f\|_p$ is an average of $|f|^p$; as $p$ grows the largest values dominate. Above, $|f|\le\|f\|_\infty$ costs nothing; below, the set where $|f|$ is close to its essential supremum has positive
measure, and $\mu(A)^{1/p}\to1$.
\writeup Let $M_\infty=\|f\|_\infty$. For $p\ge p_0$, $|f|^p\le M_\infty^{p-p_0}|f|^{p_0}$ a.e., so $\|f\|_p\le M_\infty^{1-p_0/p}\|f\|_{p_0}^{p_0/p}\to M_\infty$, and $\limsup\|f\|_p\le M_\infty$. If $M_\infty=0$ we are done. For
$0<M<M_\infty$, the set $A=\{|f|>M\}$ has $\mu(A)>0$ by definition of $M_\infty$, and $\mu(A)\le M^{-p_0}\|f\|_{p_0}^{p_0}<\infty$ (Markov). Then $\|f\|_p\ge M\mu(A)^{1/p}\to M$, so $\liminf\|f\|_p\ge M$ for every $M<M_\infty$.
The constant function $1$ on $\R$ is in $L^\infty$ and in no $L^p$ with $p<\infty$.
\end{solution}

\begin{solution}{ex:hardy}
\idea $F(x)$ is an average of $f$ over $[0,x]$, which after rescaling is the average of $f(xu)$ over $u\in[0,1]$. Each rescaled $f(x\,\cdot)$ has $L^p$ norm $u^{-1/p}\|f\|_p$, and H\"older with a weight
$u^{\pm\alpha}$ trades this growth at $u=0$ against integrability.
\writeup Substituting $s=xu$, $F(x)=\int_0^1f(xu)\,du$. Let $\alpha>0$ with $\alpha q<1$. H\"older in $u$ gives
\[F(x)^p\le\int_0^1f(xu)^pu^{\alpha p}\,du\cdot\Bigl(\int_0^1u^{-\alpha q}\,du\Bigr)^{p/q}.\]
Integrate in $x$ (Tonelli), using $\int_0^\infty f(xu)^pdx=\frac1u\|f\|_p^p$:
\[\|F\|_p^p\le\|f\|_p^p\int_0^1u^{\alpha p-1}du\cdot(1-\alpha q)^{-p/q}=\frac{\|f\|_p^p}{\alpha p\,(1-\alpha q)^{p/q}}.\]
With $\alpha=\frac1{pq}$: $\alpha p=\frac1q$ and $1-\alpha q=1-\frac1p=\frac1q$, so the bound is $q\cdot q^{p/q}=q^{1+p/q}=q^p$. Hence $\|F\|_p\le q=\frac p{p-1}$. Sharpness: for $f=x^{-1/p}\mathbf 1_{[1,N]}$,
$\|f\|_p^p=\ln N$, and for $1\le x\le N$, $F(x)=\frac1x\cdot q(x^{1/q}-1)=q(x^{-1/p}-x^{-1})$. So $\|F\|_p^p\ge q^p\int_1^N(x^{-1/p}-x^{-1})^pdx$, and since $(x^{-1/p}-x^{-1})^p=x^{-1}(1-x^{-1/q})^p$ with $(1-x^{-1/q})^p\to1$,
the integral is $\ln N-O(1)$ (the difference is $\int_1^Nx^{-1}(1-(1-x^{-1/q})^p)\,dx\le\int_1^\infty p\,x^{-1-1/q}\,dx<\infty$, using $1-(1-v)^p\le pv$). Hence $\frac{\|F\|_p^p}{\|f\|_p^p}\ge q^p(1-\frac{O(1)}{\ln N})\to q^p$.
\end{solution}

\begin{solution}{ex:jensen}
\idea Each inequality is Jensen for a well-chosen convex function and probability: $\exp$ for (a), $t^2$ for (b), and the concave $\ln$ (Jensen reversed) for (c).
\writeup (a) With the uniform probability on $\{1,\dots,n\}$, $f(i)=\ln x_i$ and $\varphi=\exp$: $\exp(\frac1n\sum\ln x_i)\le\frac1n\sum x_i$, and the left side is $(x_1\cdots x_n)^{1/n}$. (b) $\varphi(t)=t^2$ (if $\mathbb E[X^2]=\infty$ there
is nothing to prove). (c) $-\ln$ is convex on $]0,\infty[$, so $-\ln\int f\le\int(-\ln f)$, where $\int(-\ln f)$ is well defined in $]-\infty,+\infty]$ by Theorem~\ref{thm:jensen}.
\end{solution}

\begin{solution}{ex:linfinity}
\idea A separable space cannot contain uncountably many points at mutual distance $1$, since disjoint balls of radius $\frac12$ around them would each need a point of a countable dense set. In $L^p$
with $p<\infty$, the density of step functions does the counting.
\writeup (a) For $s<t$, $\|\mathbf 1_{[0,t]}-\mathbf 1_{[0,s]}\|_\infty=\|\mathbf 1_{]s,t]}\|_\infty=1$. If $\{g_n\}$ were dense, each open ball $B(\mathbf 1_{[0,t]},\frac12)$ would contain some $g_n$; these balls are disjoint, so $t\mapsto n$ would
be injective from $[0,1]$ into $\N$, which is impossible. (b) By Theorem~\ref{thm:ccdense}, step functions are dense; each step function on $[0,1]$ is approximated in $L^p$ by one with rational endpoints
and rational (or Gaussian-rational) heights, a countable family. (c) A step function, or a function of $C_c(\R)$, vanishes outside some interval $[-R,R]$. Beyond $R$ lie whole periods of $\sin$, and $\sin>1-\varepsilon$ on a set of positive measure there,
so $\|\sin-h\|_\infty\ge1-\varepsilon$ for every $\varepsilon>0$. Bounded continuous functions are thus far from $C_c$ in $L^\infty$: uniform approximation needs the behaviour at infinity to match.
\end{solution}

\section*{Sources and further reading}

\begin{itemize}
\item F.~Riesz, ``Untersuchungen \"uber Systeme integrierbarer Funktionen'', \emph{Math. Ann.} 69 (1910), 449--497. The $L^p$ spaces, H\"older and Minkowski for functions, and completeness.
\item E.~Fischer, ``Sur la convergence en moyenne'', \emph{C. R. Acad. Sci. Paris} 144 (1907), 1022--1024; F.~Riesz, ``Sur les syst\`emes orthogonaux de fonctions'', same volume, 615--619.
\item G.~H.~Hardy, J.~E.~Littlewood and G.~P\'olya, \emph{Inequalities}, 2nd ed., Cambridge University Press, 1952. H\"older, Minkowski, Jensen and Hardy's inequality, with their history.
\item G.~B.~Folland, \emph{Real Analysis}, 2nd ed., Wiley, 1999, Chapter~6 and Section~8.2.
\item E.~H.~Lieb and M.~Loss, \emph{Analysis}, 2nd ed., AMS, 2001, Chapters~2 and~4. Convolution, Young's inequality with its sharp constant, and approximation.
\item \emph{For the stories in the introduction.} L.~J.~Rogers, ``An extension of a certain theorem in inequalities'', \emph{Messenger of Math.} 17 (1888), 145--150; O.~H\"older,
``Ueber einen Mittelwerthssatz'', \emph{Nachr.\ Ges.\ Wiss.\ G\"ottingen} (1889), 38--47; J.~L.~W.~V.~Jensen, ``Sur les fonctions convexes et les in\'egalit\'es entre les valeurs
moyennes'', \emph{Acta Math.} 30 (1906), 175--193. For Jensen's career at the telephone company, see his biography in the MacTutor History of Mathematics archive
(St Andrews).
\end{itemize}
